% ===== §10 Perception: How the Field Becomes Visible =====
\section{The Field Becomes Visible: Perception and Its Contending Theories}
\label{sec:perception-field}

\epigraph{\zh{应无所住而生其心。}}{\textit{Diamond Sutra}}

\noindent
The first movement of the cycle was stated almost too gently. To perceive the field is to let appropriateness be seen, to suspend the projection of purpose and leave room for what is there to show itself. Put that way it sounds like a counsel of attentiveness, the sort of thing one nods at and forgets. Underneath the gentle formulation lie several complete and mutually incompatible theories of how perception is possible at all, each with its own mechanism, its own training, and its own metaphysical commitment, and they do not agree. A predictive brain, a mind resting in undivided awareness, and a body absorbed in a world are three different answers to the question of where seeing happens and what it would take to see better. This section sets them against one another rather than gathering them into a comfortable plurality, because the disagreements are real and because the framework of this series is not neutral among them. What looks like a single instruction, leave room for the field to appear, turns out to name at least three distinct operations, and the work of the section is to show how they stand to one another and where the present account must take sides.

\subsection{The problem: why appropriateness goes unseen}
\label{sub:perception-problem}

Begin with the phenomenon the instruction addresses. In the ordinary conduct of a shared life most of what passes goes unregistered. The table is set, the greeting returned, the cup placed, and none of it is seen, in the sense that none of it presents itself as bearing appropriateness or its lack. This is not a failure of the eyes. It is the normal condition of competent absorption in a familiar world, and it has a function: a life in which every placement of every cup demanded fresh appraisal would be unlivable. The cost of the function is that the field becomes invisible exactly where it is most lived in. Appropriateness goes unseen not because it is hidden but because seeing it requires the interruption of a fluency that exists precisely in order not to be interrupted.

So the instruction to leave room is an instruction to interrupt something, and the disagreements among the traditions can be located with some precision once we ask what, exactly, is to be interrupted. One tradition says the thing to be loosened is the grip of the prior, the top-down expectation that flattens the field into what was predicted. Another says the thing to be loosened is conceptual division itself, the movement by which awareness, the instant it arises, is overlaid with names and with the split into a seer and a seen. A third says the thing to be loosened is the absorption of the thing into its use, the readiness-to-hand in which the cup disappears into the drinking. These are not three names for one operation. They cut at different joints, they prescribe different trainings, and whether they are aspects of a single thing seen from three sides is a question to be argued rather than assumed. The remainder of the section states each in turn and then asks where they conflict.

To keep the contest concrete, it helps to fix a single ordinary case and to ask of each tradition what it would say about the same small event, since the disagreements are easiest to feel when they are disagreements about one thing rather than three.

\begin{casebox}[A morning, and a sentence that was not foreseen]
Two people share a morning that has its shape. The coffee is made in the order it is always made, the table laid as it is always laid, and into this familiar shape the one says something the other did not expect, a remark slightly off the rhythm of the mornings before it, carrying a small weight that the routine did not predict. The remark is a contingency, a thing that need not have been said and was. It can pass entirely unregistered, absorbed into the morning's fluency as though it had been one more expected motion; or it can be seen, registered as bearing something, a slight shift in the field that asks to be met. The whole question of perceiving the field is the question of what divides the morning in which the remark is absorbed without trace from the morning in which it is seen, and each of the traditions that follows gives that division a different name.
\end{casebox}

\subsection{The predictive route}
\label{sub:perception-predictive}

On the first account, perception is not the reception of a given but the brain's active construction of its best hypothesis about the causes of its input, continually corrected by the error between what was predicted and what arrived~\citep{clark_whatever,friston_free}. Seeing is inference. To see appropriateness in the field is, in this language, to keep the weight assigned to incoming evidence from being overwhelmed by the weight assigned to the prior, so that the small deviation, the unplanned thing, the way the light fell or the way the answer came, is not silently explained away by what was expected. The relevant variable is the balance of precision: how much the system trusts its prediction relative to its evidence. A strong prior, heavily weighted, flattens the field, because it can account for almost anything in advance and so registers almost nothing as surprising. Training, on this account, is the modulation of that balance, a learned refusal to let expectation suppress the senses, and it has the great advantage of being measurable, since the system's responsiveness to deviation leaves traces that can be looked for.

This route is powerful and it is the one the implementation section will lean on, but its metaphysical commitment must be marked, because it is the commitment that will later force a choice. The predictive account is naturalistic and representational. Perception is inference performed by an internal model, and the model is housed in a single cranium, building its estimate of an external world from sensory evidence. Seeing the field, on this account, is one organism's improved estimate of its surroundings. That is a coherent and fertile picture, and it sits in tension with the claim of relational aesthetics that the perception in question has its proper site in the coupling rather than in either pole alone. A theory on which seeing is what one brain does to its input cannot, without modification, be the whole theory of a seeing whose object is an event between. The tension is not yet a refutation. It is a debt, and the section records it now in order to discharge it later.

On this account the morning's remark is absorbed or seen according to the precision the hearer's system assigns. Where the prior is heavily weighted, the morning has been so thoroughly predicted that the remark is explained away before it can register, folded into the expected shape of every prior morning, its small deviation absorbed by a model confident enough to account for it in advance; the hearer has, in the literal sense the account intends, predicted the remark out of existence. Where the weight on evidence is kept high enough, the deviation survives, the remark arrives as a genuine surprise that the model did not anticipate, and the prediction error it raises is the signal that something in the field has shifted. The training the account recommends is the learned lowering of the prior's weight, the cultivated refusal to let the confident model of the morning suppress the evidence of this morning, and its great strength is that this lowering should be visible, a measurable change in how much the system responds to the deviation it once would have absorbed. What the account cannot say, and does not pretend to say, is anything about the two who share the morning as two; the remark is, for it, an input to one hearer's model, and the seeing of it is that model's improved estimate, which is precisely the limitation that will force the account to its proper level later.

\subsection{The Buddhist route}
\label{sub:perception-buddhist}

The second account stands in the sharpest opposition to the first, and it is the one most consonant with the roots this series has drawn on. Its diagnosis is that appropriateness goes unseen because perception, the instant it arises, is covered over by conceptual division: by names, by the sorting of what appears into kinds already held, and above all by the split into a perceiving subject set over against a perceived object. The single-practice absorption named in the early sutras and reread by the Chan tradition, \zh{一行三昧} (\textit{yixing sanmei}), is not a withdrawal from the ordinary into a special state but a sustained, undivided directness maintained through walking, standing, sitting, and lying, in which moment after moment no discrimination is raised. Alongside it the contemplative analysis of \zh{止观} and the Yogacara category of direct perception, \zh{现量} (\textit{pratyaksa}), the immediate apprehension that precedes the imposition of name and concept, point together at one thing: what obscures the field is not a deficit of representation but the very activity of representing, the reflex by which bare awareness is instantly furnished with a subject, an object, and a label binding them.

To see appropriateness, on this account, is therefore not to represent the field more accurately but to relax the structure of representation as such, to let what appears appear before the split that would set a viewer over against a viewed. The training is neither the modulation of a parameter nor the disruption of habit but a kind of keeping: to let awareness rest, undivided, so that when discrimination arises it is noticed and not followed. The metaphysical commitment is non-dual and anti-representational, and it is for this reason that the Buddhist route cannot simply be folded into the predictive one. Where predictive processing says perception is the inference of a subject about an object, this tradition says that the subject-object structure is itself the obscuration, and that the deepest seeing is the one in which that structure has gone quiet. Held together, the two accounts do not describe the same operation at different levels. They make incompatible claims about whether the seer-seen structure is the apparatus of perception or the veil over it, and the section will have to say how that incompatibility is to be borne.

The Chan rereading of the single-practice absorption sharpens this and removes a likely misunderstanding. When the Platform Sutra glosses the practice, it refuses the picture of a special seated state cut off from ordinary life, the picture of one who would sit fixed and empty and call that the practice; the single practice is rather the maintaining of a direct and undivided heart in all the ordinary postures, walking and standing and sitting and lying, throughout every activity of the day~\citep{platform_sutra}. The point is exactly that the practice belongs to the morning and the coffee and the laid table and not to a withdrawal from them, which is what makes it a candidate for an aesthetic education of the everyday rather than an escape from the everyday into contemplation. To rest undivided is not to cease the morning's activity but to conduct it without raising, at each step, the grid of name and verdict and the split of a self against its world; awareness stays direct, and when the dividing reflex arises it is seen and not followed. This is a discipline performed in the midst of the daily and not beside it, and that is the feature the section will need when it asks how any of these trainings could be available in a life with no leisure to withdraw.

Returned to the morning, the Buddhist account locates the obscuration not in a heavily weighted prior but in the reflex that, the instant the remark arrives, overlays it with a verdict and a position: the hearer's instantaneous sorting of the remark into a kind already held, the silent placing of it as the sort of thing she says, the constitution of a hearer set over against a speaker who is thereby fixed as an object understood. On this account the remark goes unseen not because it was predicted but because it was immediately known, captured by a discrimination so fast it feels like perception itself, and what is needed is not the lowering of a weight but the quieting of the dividing reflex, so that the remark may be met before it has been placed, in a directness that has not yet split into a knower appraising a known. The difference from the predictive reading is not a matter of emphasis. The predictive account would have the hearer weight the evidence of the remark more heavily; the Buddhist account would have the hearer stop performing, for a moment, the very act of knowing-as that the predictive account takes perception to be.

\subsection{The phenomenological route}
\label{sub:perception-phenomenological}

The third account lies between the other two and shares the anti-representational instinct of the second without its renunciation of the subject. In the analysis of everyday existence the world is encountered first as ready-to-hand: things are absorbed into their uses, the cup into the drinking, the door into the passing-through, and in this absorption they precisely do not appear as objects with properties. What is most used is least seen. It is when the familiar is broken, when the tool fails or the routine is interrupted, that the thing comes forward and is lit up as what it was, so that appropriateness is most often disclosed at the very point where ordinary fluency is broken off. To this the phenomenology of perception adds that seeing is never in the first place the operation of a mind upon a representation; it is the bodily involvement of a perceiver already in the world, prior to the division into subject and object, the flesh of the perceiver and the flesh of the world caught in a single reversible fold~\citep{merleau_phenomenology}.

The training that follows from this account is neither the inward keeping of the Buddhist route nor the parameter-tuning of the predictive one. It is the cultivation of a wakefulness to fallen fluency, a practice of letting the familiar grow strange so that what has sunk into use can come forward again and be seen. Art does this, and so does the simple discipline of stopping before a thing one uses every day and looking at it as though for the first time. The metaphysical commitment is anti-representational but retains the perceiver's embodiment and being-in-the-world, and this is why the route sits between the others. It refuses, with the Buddhist account, the picture of a mind representing an external world; it retains, against that account, a perceiving body that is genuinely a subject, though a subject already mingled with what it perceives. The figure of the flesh, the reversible belonging of perceiver and perceived, is very nearly the relational coupling this paper needs, and it is no accident that of the three routes this is the one from which the bridge to relational aesthetics will be shortest.

On this third account the morning's remark goes unseen for a reason different again from the heavy prior and the dividing reflex: it goes unseen because the morning has fallen into use, because the shared breakfast has become a smooth machinery of motions in which each person and each act is absorbed into its function, the other absorbed into her role in the morning as the cup is absorbed into the drinking. In a morning lived wholly as ready-to-hand the remark is not absorbed by a confident prediction nor placed by a fast verdict but simply passed over, because nothing in the fluent machinery comes forward to be seen at all, the other least of all, she having sunk into her use in the morning's smooth running. What lets the remark appear is the small breaking of the fluency, the catch in the machinery where the routine stumbles and the other comes forward, no longer absorbed into her function but standing out as herself, present in a way the fluent morning had kept her from being. The training, accordingly, is the cultivation of a wakefulness that does not wait for the fluency to break of itself but lets the familiar morning grow strange enough that the other can come forward within it, a practised refusal to let the shared life sink so wholly into use that the one who shares it disappears into her role. The bodily accent matters here: the seeing is not a mind's better estimate but the perceiver's whole embodied involvement in the morning catching, for a moment, on the other as a presence and not a function, and that catching is already half of the reversible fold in which the relational coupling will be found.

\subsection{A Daoist second voice}
\label{sub:perception-daoist}

Before turning to the conflicts, one further voice should be admitted, close to the Buddhist route but not identical with it, and belonging to the same roots. The fasting of the heart in the Zhuangzi, the injunction to attain utmost emptiness and hold to stillness in the Laozi~\citep{zhuangzi,laozi_daodejing}, describe a receptivity that meets things by first emptying the self of the formed mind that would meet them with a verdict already prepared. Where the Buddhist route emphasises the breaking of discrimination, this Daoist voice emphasises the emptying of the self so that the thing may present itself of its own accord, \zh{虚而待物}, vacant and awaiting the thing. The difference is one of accent rather than of doctrine, but the accent matters: the first guards against the imposition of concept, the second against the imposition of the self's prior intent, and the two together mark out, from the Eastern side, both of the impositions that the instruction to leave room is trying to suspend. The Daoist voice will return when the practical training is described, because emptying-to-await is, of the four, the figure that translates most directly into something one can do before an ordinary act.

\subsection{A bridge: perceptual expertise}
\label{sub:perception-expertise}

Between the predictive route and the phenomenological one lies a body of empirical work on perceptual expertise that the section should admit, because it occupies the ground the two share and shows the disagreement between them to be narrower, in one respect, than their metaphysical commitments suggest. The study of expert perception, of how the trained radiologist sees the anomaly the novice misses, how the experienced birder identifies in a glance what the beginner cannot resolve, how the chess master perceives the position as a structure where the amateur sees scattered pieces, establishes that perception itself, and not merely judgement downstream of it, is transformed by training. The expert does not reason faster from the same percept; she perceives differently, the relevant structure standing out for her where it was invisible before, and this transformation is the lowering of a threshold, the acquisition of a discrimination that has become perceptual rather than inferential. This is the empirical correlate of what the paper calls the raising of resolution, and it sits between the predictive and phenomenological routes because it can be described in either idiom: as a learned change in the precision-weighting that lets the expert's system register what the novice's predicts away, or as the coming-forward of a structure that for the expert is no longer absorbed into an undifferentiated field. The two idioms describe, here, one well-documented phenomenon, the perceptual transformation wrought by training, and the convergence of the predictive and phenomenological vocabularies upon it suggests that their disagreement, real at the level of ultimate commitment, leaves a wide middle ground where they describe the same thing in different words.

The relevance to the cultivation of the field is direct, and it answers a doubt that might otherwise stand. One might doubt whether the perception of a field's appropriateness can really be trained, whether it is the kind of thing that improves with exercise or merely a sensitivity one has or lacks. The literature on perceptual expertise removes the doubt by establishing, across many domains, that perception of exactly this kind, the discrimination of a structure that no rule specifies, recognised in the concrete and not derived, is trainable, that exposure and exercise transform it, that the threshold for registering the relevant structure lowers with practice. If the perception of a chess position or a radiograph can be trained from coarse to fine, so, the analogy suggests, can the perception of a relational field, which is a discrimination of the same kind, the recognition of a structure in the concrete that no rule captures. The empirical work does not prove that the field's appropriateness is trainable, since the analogy is not a demonstration, but it establishes that perception of the relevant kind is trainable in general, which removes the doubt that the field might be the one domain where such training is impossible, and lends the paper's central claim, that the perception of appropriateness can be cultivated, the support of a wide body of evidence that perception of its type can be.

\subsection{Three real conflicts}
\label{sub:perception-conflicts}

It would be easy, and false, to arrange these accounts as complementary emphases on a single underlying process. They conflict, and the conflicts are worth stating sharply, because the framework of this series resolves them by taking sides rather than by blending.

The first conflict concerns the site of perception. The predictive route places it inside the cranium, as the inference of an internal model. The phenomenological route places it between, in the bodily fold where perceiver and perceived belong to one another before they are divided. The Buddhist route places it nowhere that the question presupposes, since the very framing of an inside that perceives an outside is, on its account, the obscuration to be undone. These three answers cannot all be true. A seeing that is one brain's estimate of its surroundings is not a seeing that occurs in the reversible fold of flesh, and neither is a seeing in which the inside-outside structure has gone quiet. The framework of relational aesthetics has already committed itself on this point: the perception that matters has its site in the coupling, not in either pole. That commitment is incompatible with taking the predictive account as the whole story, and it leans hard toward the phenomenological and the Buddhist side. The predictive route, on pain of contradicting the concept the paper has built, must be retained as something other than an account of where seeing is. It is retained, as the implementation section will make precise, as the description of an observable trace, not as the ontology of the act.

The second conflict concerns what is to be suspended. The phenomenological route suspends the usefulness of the ready-to-hand, so that the thing may come forward. The Buddhist route suspends conceptual division, so that the seer-seen split may go quiet. The predictive route suspends nothing in the strict sense; it rebalances, lowering the weight of the prior so that evidence is not overwhelmed. These are different operations with different objects, and a paper that says only leave room for the field has not yet said which suspension it means. The honest answer is that the instruction names all three, pitched at different grains, and that the burden of the present account is to show they are not three unrelated disciplines but three depths of one descent: the rebalancing of the prior is the coarsest and most measurable, the de-automation of use is intermediate, and the quieting of division is the deepest and least capturable. This is a claim, not a definition, and it is argued rather than stipulated. Its argument is that each deeper suspension includes the shallower as a special case: one cannot quiet the seer-seen division while leaving the use-absorption intact, and one cannot de-automate use while leaving a heavily weighted prior to flatten the field in advance. The inclusion runs one way and not the other, which is why the three can be ordered as depths rather than merely juxtaposed as alternatives.

\subsection{The hardest step: non-division grounds discrimination}
\label{sub:perception-hardest}

The third conflict is the dangerous one, and meeting it is the labour the section exists for. The Buddhist route demands the relinquishing of discrimination; this paper demands that the perceiver \emph{discriminate}, finely, whether the field is accumulating holonomy or spinning in place, what here is appropriate and what is not. On its face the demand for non-discrimination cancels the discriminating power the account needs. If to see well is to let all division go quiet, how can to see well also be to register, with great delicacy, the difference between a coupling that gains and one that merely turns? The objection is not a quibble. It threatens to make the Buddhist route and the framework of this series simply inconsistent, so that one of them must be given up.

The resolution turns on a distinction internal to the Buddhist analysis itself, and it is the same shape as the distinction this series has drawn elsewhere between two kinds of repetition and two kinds of trust. What the practice dissolves is the deluded discrimination, \zh{妄分别}, the conceptual, subject-against-object sorting that overlays bare awareness with a grid of prior verdicts. It does not dissolve, and in the developed accounts it positively makes possible, the subtle discernment that the tradition calls the wisdom of wondrous discernment, \zh{妙观察智}, the differentiating knowledge that belongs to the awareness recovered after the deluded grid has gone quiet. The order is the whole point. First the deluded division is let go; only then does appropriateness show itself as it is, undistorted by the prior that would have pre-sorted it; and the fine discrimination of good accumulation from bad is exercised upon what is then disclosed.

\claim{Non-division grounds discrimination rather than abolishing it. What contemplative practice quiets is the deluded, conceptual, subject-against-object division that pre-sorts the field; what it makes possible is the subtle discernment of appropriateness, which can be exercised only upon a field that has first been let appear undistorted. The perception of holonomy that this account requires is not the survival of ordinary discrimination through the practice but its refinement on the far side of the practice. Letting go of the deluded grid is the condition of the finer seeing, not its casualty.}

A second and harder form of the objection presses past this, and it must be met or the resolution will seem too easy. Grant the distinction between deluded and subtle discrimination; still, the framework asks the perceiver not merely to discern finely but to \emph{care} whether the field gains or spins, to be invested in the difference, to want the coupling to accumulate, and this caring, the objection runs, is exactly the grasping, the abiding, the attachment to an outcome that the contemplative traditions name as the deepest obscuration of all. A perception that wants the field to gain is a perception bent by desire, and a desire-bent perception is precisely what cannot see what is, since it sees what it wishes. On this stronger form, it is not discrimination that the practice and the framework fight over but investment, and the framework's investment in accumulation seems to reintroduce, at the level of caring, the very grasping the practice exists to release.

The answer is given by the line that stands at the head of this section, and it is the most exact possible statement of the resolution: to give rise to the heart that abides nowhere. The teaching does not say to give rise to no heart, to care about nothing, to perceive without investment; it says to give rise to the heart, to let the caring arise, while abiding nowhere, without fixing it upon a possessed outcome clung to against loss. The caring the framework requires is of this kind, or can be cultivated to be: an investment in the field's gaining that does not grasp at a particular result, a wanting the coupling to accumulate that does not seize the accumulation as a thing to be had and defended. This is the difference between desiring that the field gain and demanding that it gain in a specified way by a specified time, between a care that remains open to what the field actually does and an attachment that has already decided what it must do and sees only its own demand. The non-abiding heart cares and does not grasp, perceives the field's direction with full investment and without the seizure that would bend the perceiving, and it is exactly this heart, caring yet unabiding, that the framework's perceiver must cultivate. Far from reintroducing the grasping, the resolution shows the framework's investment to be the very thing the teaching describes, a heart that arises and abides nowhere, and the strongest form of the objection inverts like the first: the contemplative tradition does not forbid the caring the framework needs but specifies its non-grasping form.

This second inversion is the same movement as the first and belongs with it. The first showed that letting go of deluded division founds the subtle discernment; the second shows that letting go of grasping investment founds the open caring; and both are the one counsel that runs through every step of this paper, to generate without possessing, brought to the conditions of perception. To see the field well is to discriminate finely without the deluded grid and to care fully without the grasping seizure, and these two are not separate achievements but one, the single bearing of a perception that meets the field invested and unattached, discerning and undivided, which is what the cultivation of the beauty of perception finally trains.

If this is right, then the most threatening of the three conflicts inverts into the deepest support. The Buddhist route does not take away the discriminating power the framework needs; it grounds it, by clearing the prior that would otherwise have done the discriminating in advance and badly. This is the same movement that the paper performs at each of the four steps, the movement in which a relinquishing of grasp turns out to found rather than forbid a generation. Here it is the relinquishing of conceptual seizure that founds the finer sight. The shape will recur, at reception as the mirror that meets without storing, at sharing as the union that preserves rather than dissolves difference, at reproduction as the accumulation that does not hoard. Each is one face of the same ancient counsel, to generate without possessing and to act without holding, read into the conditions of perception, of welcome, of the between, and of time.

\subsection{A layered synthesis}
\label{sub:perception-synthesis}

The three routes can now be placed, each at the level where it speaks without overreaching, and the placement uses no equals sign, because none of the relations among them is an identity and none survives being asserted as one. The phenomenological route is taken at the level of ontology: it gives the best available picture of what perceiving the field is, a bodily belonging of perceiver and perceived prior to their division, and it is from this picture that the bridge to the coupling of relational aesthetics is shortest, since the reversible fold of flesh is already a between. The Buddhist route is taken at the level of cultivation and of metaphysical guard: it gives the discipline by which the prior is quieted and the field let appear, and it holds the anti-representational line that keeps the whole account from collapsing back into one brain estimating an outside, while its internal distinction between deluded and subtle discrimination secures the discriminating power the framework cannot do without. The predictive route is taken at the level of implementation: it corresponds to the observable rebalancing of precision, and an account on which seeing the field leaves a measurable signature in a system's responsiveness to deviation may be read alongside the ontological and cultivational accounts without claiming to be either. The Daoist voice runs through all three as the figure of emptying-to-await, the posture that the practical training will name.

What holds the synthesis together is not a doctrine common to the three but the shape of their cooperation, which is itself relational. The traditions are incommensurable in their commitments and are not made commensurable here; they are each given their proper place and asked to do only what they can, which is the seeking of common ground while the differences are kept, the method this series has named and practised. The synthesis is therefore not a higher theory that absorbs the three. It is an arrangement in which each is preserved as itself and none is rendered mere material for another, and that, by the paper's own lights, is the only kind of synthesis that would not contradict what the paper is about.

One might object that the assignment of levels is convenient rather than principled, that calling the phenomenological account ontological, the Buddhist account cultivational, and the predictive account implementational is a way of avoiding their conflict by fiat rather than resolving it. The objection deserves an answer, and the answer is that the assignment follows from what each account is best at and worst at, not from the convenience of the paper. The predictive account is assigned the implementation level because it is there that it is strongest and there that its representational commitment does no harm: as an account of the observable machinery by which a brain participates in perception, it is precise, measurable, and fertile, and its picture of an internal model housed in a cranium is exactly right as a description of the implementing mechanism, becoming false only when it is promoted to the ontology of the act, which is why it is kept below the ontology. The phenomenological account is assigned the ontological level because its picture of the reversible fold is the one that correctly describes what perceiving the field is, a bodily belonging prior to the subject-object split, and because that picture is the one continuous with the relational ontology the series requires; it is weakest as an implementation account, having little to say about mechanism, which is why it is kept above the implementation. The Buddhist account is assigned the cultivational level and the office of metaphysical guard because its strength is the discipline it describes and the anti-representational line it holds, and because its claims about the deepest structure of awareness, while they cannot be made the paper's official ontology without committing the paper to more than it can defend, must be honored as the guard that keeps the ontology from collapsing into representation. The levels are thus assigned by the fit between each account's strengths and the work of each level, and the assignment is principled in exactly this sense, that each account is placed where it does its best work and kept from the level where its commitments would mislead.

\subsection{The reach and the defeat of each route}
\label{sub:perception-ledger}

The layered arrangement just defended can seem, stated as an assignment of levels, to settle the contest too smoothly, as though each tradition were simply handed a province and sent away content. It is worth being blunt, then, about what each route explains that the others cannot, and about what defeats each taken alone, because the arrangement earns its authority only if it is built on a frank accounting of strengths and failures rather than on a diplomatic distribution of honours. The accounting also discharges a debt the paper owes its own commitments: holding, as it does, to a dialectical and materialist refusal to crown any theory as the final rationality, it must show in each case why the theory cannot be the whole, and not merely assert that it occupies a level.

The predictive route explains what neither of the others can, the graded, trainable, measurable character of perceptual sensitivity, the fact that the registration of the field improves with exercise in ways that leave physical traces and admit of degrees. It is the only one of the three that connects the perception of the field to a mechanism that could be studied, perturbed, and measured, and to discard it would be to sever the account from the empirical altogether. What defeats it, taken alone, is the betweenness of the object it would explain. A theory on which perceiving is one brain's inference about its input has no resources for a perceiving whose object is a coupling between two subjects, and when it is pressed to account for the shared perception that relational aesthetics requires, it must either deny that the betweenness is real, which contradicts the concept the paper has built, or smuggle in a relational structure its own terms do not contain. Its representationalism, which is its strength as a mechanics, is its defeat as an ontology, because the relational event it would have to represent is not in the single head where it locates all representing. This is the precise sense in which the materialist commitment bites: the predictive route is retained for the real work it does at the level of mechanism and refused at the level where its representationalism would dissolve the relational reality into a fact about one nervous system.

The phenomenological route explains what the predictive route cannot, the lived structure of perceiving as a bodily belonging to a world prior to the division of subject and object, and it is from this structure that the bridge to the coupling is built, since the reversible fold of perceiver and perceived is already a between. It is strongest exactly where the predictive route is weakest, at the ontology of the act, and its account of fallen fluency and the coming-forward of the broken-from gives the phenomenology of perception a precision the other routes lack. What defeats it, taken alone, is its silence about mechanism and its difficulty with cultivation. It describes beautifully what perceiving is and says little about how a perception is trained from coarse to fine, offering no account of the graded improvement that the predictive route and the literature on expertise document, and a philosophy of aesthetic education cannot rest on an account that illuminates the structure of perception while leaving dark the very trainability that makes education possible. Its strength is description and its defeat is the silence of description before the question of cultivation.

The Buddhist route explains what neither of the others addresses at all, the obscuration that perception inflicts upon itself, the way the registering of the field is veiled not by a deficit but by the very reflex of representation, and it alone offers a discipline aimed at that obscuration, together with the distinction between deluded and subtle discrimination on which the paper's hardest step turned. What defeats it, taken as the whole truth rather than as a discipline and a guard, is its tendency, pressed to its metaphysical limit, to dissolve the very field whose appropriateness the paper needs perceived. A doctrine on which the subject-object structure is wholly an obscuration to be quieted threatens, at its term, to quiet also the discrimination of the field's generative direction that the cultivation requires, and although the paper argued that non-division founds rather than abolishes the subtle discernment, that argument holds the Buddhist route at a determinate point and does not follow it all the way to the dissolution of all distinction. The route is honoured as cultivation and as the guard against representationalism, and it is not adopted as a final metaphysics, because a final non-dualism would take with it the discriminated field the paper cannot do without. Here the materialist commitment bites from the other side: the contemplative route is refused the status of ultimate ontology not because its discipline is doubted but because its metaphysical terminus would dissolve the differentiated material field in which alone the good cycle is built.

\subsection{The materialist resolution and its philosophical position}
\label{sub:perception-materialist}

These three accountings, set together, do not yield a synthesis that absorbs the three into a higher theory, and the refusal of such a synthesis is not a failure of nerve but the paper's considered position, which is worth naming for what it is. The position is dialectical and materialist. It is materialist in that it takes the field, the coupling, the bodily belonging, and the trainable mechanism to be features of a real world that perception is of, not constructions of a transcendental subject nor moments of an absolute idea unfolding, and it treats the perception of the field as a real material process with a real material object, the relational field whose accumulation is as real as any other material accumulation. It is dialectical in that it refuses to let any one of the three accounts stand as the whole truth, holding instead that each grasps a real aspect, the mechanism, the lived structure, the self-obscuration, and that the truth is the articulated relation of these aspects rather than any one of them, an articulated relation in which each is determinately placed and determinately limited. This is not eclecticism, the indifferent borrowing of whatever is convenient from each, because the placement is principled, each account assigned the level where it grasps a real aspect and refused the level where its commitments would distort, and the principle of the placement is the materialist insistence that the relational field is real and the dialectical insistence that no single account of a real thing is its whole truth.

The position can be located against the alternatives it refuses. It refuses the reductive naturalism that would make the predictive account the whole, dissolving the relational field into a fact about nervous systems, because that reduction loses the betweenness that is the field's reality. It refuses the transcendental idealism that would make the structure of the subject's faculties the ground of the perceived, because that idealism loses the materiality of the field, making the coupling a construction of the perceiving subject rather than a real relation the subject enters. And it refuses the contemplative absolutism that would make non-division the final truth, dissolving the differentiated field into an undifferentiated awareness, because that dissolution loses the discriminated direction the cultivation requires. Against all three the paper holds that the relational field is a real material structure, perceived by a real material process, graspable in its mechanism and its lived form and its liability to self-obscuration, and reducible to none of these. This is the materialist dialectic applied to the theory of perception, and it is the same refusal of any theory's claim to be the final rationality that governs the paper throughout, here issuing not in a synthesis of the three routes but in their determinate articulation, each doing the real work it can and none permitted to be the whole.
